% GENERATED FILE. Edit canonical lessons, then run npm run generate:books.
% canonical-source-hash: fnv1a64:a94613244376b116
% canonical-lessons: LA-C390-condicio, LA-C390-sin, LA-C390-dummodo, LA-C390-modo, LA-C390-quamvis

\chapter{On One Condition}
\label{ch:la-on-one-condition}
\hlchaptermodality{390}

\begin{chapteropening}
\textbf{By the end of this chapter:} \emph{I can set a condition, say but if, provided that and only if, and concede however much.}

Complete the last lesson of chapter 390: I can set a condition, say but if, provided that and only if, and concede however much.
\end{chapteropening}

\section[condiciō, condiciōnis]{condiciō, condiciōnis — a condition, terms; a situation}

\label{lesson:LA-C390-condicio}



\begin{quote}
Before the new one: say the Latin for a load, then the Latin for compliance.
\end{quote}

\subsection*{You'll want to know: condiciō, condiciōnis}
\textbf{condiciō, condiciōnis} — \textquotedblleft{}a condition, terms; a situation\textquotedblright{}.

\begin{quote}
Hāc condiciōne veniam, ut tū quoque veniās.
\end{quote}

\begin{quote}
\emph{I will come on this condition, that you come too.}
\end{quote}

\emph{condiciō, condiciōnis} f. \emph{Hāc condiciōne ut} + subjunctive: on condition that. It is also a person's situation in life. English \emph{condition}.

\begin{cousinweb}
\emph{condiciō} is built on \emph{condīcō} \textquotedblleft{}to agree on, to settle\textquotedblright{} (EtymDB; Lewis \& Short), from \emph{con-} \textquotedblleft{}together\textquotedblright{} and \emph{dīcō} \textquotedblleft{}to say\textquotedblright{}: what two sides say together. English \emph{condition} comes from it.
\end{cousinweb}

\begin{grammarlens}[title={What you've built}]
One more word for telling, explaining and describing at B1.
\end{grammarlens}

\subsection*{Your turn}
\begin{itemize}
\raggedright
  \item \emph{Say it:} \emph{condiciō, condiciōnis}
  \item \emph{Say it:} \emph{condiciō, condiciōnis}, once more
  \item \emph{Say it:} a sentence of your own with \emph{condiciō, condiciōnis}
  \item \emph{Recall:} say the Latin for a load, then the Latin for compliance, then say \emph{condiciō, condiciōnis} again
\end{itemize}

\begin{tcolorbox}[breakable,colback=teal!4,colframe=teal!35!black,title={Before you move on}]
What does condiciō, condiciōnis mean? (\textquotedblleft{}A condition, terms; a situation\textquotedblright{}.) Say it once more.
\end{tcolorbox}
\section[sīn]{sīn — but if}

\label{lesson:LA-C390-sin}



\begin{quote}
Before the new one: say the Latin for compliance, then the Latin for a condition.
\end{quote}

\subsection*{You'll want to know: sīn}
\textbf{sīn} — \textquotedblleft{}but if\textquotedblright{}.

\begin{quote}
Sī venīs, gaudeō; sīn minus, domī maneō.
\end{quote}

\begin{quote}
\emph{If you come, I am glad; but if not, I stay at home.}
\end{quote}

\emph{Sīn} opens the second of two conditions, after a \emph{sī} clause: if … but if. \emph{Sīn minus} or \emph{sīn aliter}, but if not.

\begin{cousinweb}
\emph{sīn} is built from \emph{sī} \textquotedblleft{}if\textquotedblright{} and the particle \emph{-ne} (EtymDB; Lewis \& Short): if, on the other hand.
\end{cousinweb}

\begin{grammarlens}[title={What you've built}]
One more word for telling, explaining and describing at B1.
\end{grammarlens}

\subsection*{Your turn}
\begin{itemize}
\raggedright
  \item \emph{Say it:} \emph{sīn}
  \item \emph{Say it:} \emph{sīn}, once more
  \item \emph{Say it:} a sentence of your own with \emph{sīn}
  \item \emph{Recall:} say the Latin for compliance, then the Latin for a condition, then say \emph{sīn} again
\end{itemize}

\begin{tcolorbox}[breakable,colback=teal!4,colframe=teal!35!black,title={Before you move on}]
What does sīn mean? (\textquotedblleft{}But if\textquotedblright{}.) Say it once more.
\end{tcolorbox}
\section[dummodo]{dummodo — provided that, so long as}

\label{lesson:LA-C390-dummodo}



\begin{quote}
Before the new one: say the Latin for a condition, then the Latin for but if.
\end{quote}

\subsection*{You'll want to know: dummodo}
\textbf{dummodo} — \textquotedblleft{}provided that, so long as\textquotedblright{}.

\begin{quote}
Omnia tibi dabō, dummodo vērum dīcās.
\end{quote}

\begin{quote}
\emph{I will give you everything, provided that you tell the truth.}
\end{quote}

\emph{Dummodo} takes the subjunctive; \emph{dummodo nē}, provided that not. It is \emph{dum}, while, with an adverb meaning only.

\begin{cousinweb}
\emph{dummodo} is built from \emph{dum} \textquotedblleft{}while, so long as\textquotedblright{} and \emph{modo} \textquotedblleft{}only\textquotedblright{} (Lewis \& Short): only so long as.
\end{cousinweb}

\begin{grammarlens}[title={What you've built}]
One more word for telling, explaining and describing at B1.
\end{grammarlens}

\subsection*{Your turn}
\begin{itemize}
\raggedright
  \item \emph{Say it:} \emph{dummodo}
  \item \emph{Say it:} \emph{dummodo}, once more
  \item \emph{Say it:} a sentence of your own with \emph{dummodo}
  \item \emph{Recall:} say the Latin for a condition, then the Latin for but if, then say \emph{dummodo} again
\end{itemize}

\begin{tcolorbox}[breakable,colback=teal!4,colframe=teal!35!black,title={Before you move on}]
What does dummodo mean? (\textquotedblleft{}Provided that, so long as\textquotedblright{}.) Say it once more.
\end{tcolorbox}
\section[modo]{modo — only; a moment ago; provided that (with the subjunctive)}

\label{lesson:LA-C390-modo}



\begin{quote}
Before the new one: say the Latin for but if, then the Latin for provided that.
\end{quote}

\subsection*{You'll want to know: modo}
\textbf{modo} — \textquotedblleft{}only; a moment ago; provided that (with the subjunctive)\textquotedblright{}.

\begin{quote}
Veniam, modo nē pluat.
\end{quote}

\begin{quote}
\emph{I will come, provided it does not rain.}
\end{quote}

\emph{Modo} is the ablative of \emph{modus}, a measure: only. \emph{Modo} + subjunctive, provided that; \emph{nōn modo … sed etiam}, not only … but also.

\begin{cousinweb}
\emph{modo} is the old ablative of \emph{modus} \textquotedblleft{}a measure, a limit\textquotedblright{} (Lewis \& Short): within measure, so only.
\end{cousinweb}

\begin{grammarlens}[title={What you've built}]
One more word for telling, explaining and describing at B1.
\end{grammarlens}

\subsection*{Your turn}
\begin{itemize}
\raggedright
  \item \emph{Say it:} \emph{modo}
  \item \emph{Say it:} \emph{modo}, once more
  \item \emph{Say it:} a sentence of your own with \emph{modo}
  \item \emph{Recall:} say the Latin for but if, then the Latin for provided that, then say \emph{modo} again
\end{itemize}

\begin{tcolorbox}[breakable,colback=teal!4,colframe=teal!35!black,title={Before you move on}]
What does modo mean? (\textquotedblleft{}Only; a moment ago; provided that (with the subjunctive)\textquotedblright{}.) Say it once more.
\end{tcolorbox}
\section[quamvīs]{quamvīs — however much; although (with the subjunctive)}

\label{lesson:LA-C390-quamvis}



\begin{quote}
Before the new one: say the Latin for provided that, then the Latin for only.
\end{quote}

\subsection*{You'll want to know: quamvīs}
\textbf{quamvīs} — \textquotedblleft{}however much; although (with the subjunctive)\textquotedblright{}.

\begin{quote}
Quamvīs dīves sit, nōn est fēlīx.
\end{quote}

\begin{quote}
\emph{However rich he may be, he is not happy.}
\end{quote}

\emph{Quamvīs} = \emph{quam} + \emph{vīs}, as much as you like; in classical prose it takes the subjunctive, while \emph{quamquam}, although, takes the indicative.

\begin{cousinweb}
\emph{quamvīs} is built from \emph{quam} \textquotedblleft{}as\textquotedblright{} and \emph{vīs} \textquotedblleft{}you wish\textquotedblright{}, from \emph{volō} (Lewis \& Short): as much as you like.
\end{cousinweb}

\begin{grammarlens}[title={What you've built}]
One more word for telling, explaining and describing at B1.
\end{grammarlens}

\subsection*{Your turn}
\begin{itemize}
\raggedright
  \item \emph{Say it:} \emph{quamvīs}
  \item \emph{Say it:} \emph{quamvīs}, once more
  \item \emph{Say it:} a sentence of your own with \emph{quamvīs}
  \item \emph{Recall:} say the Latin for provided that, then the Latin for only, then say \emph{quamvīs} again
\end{itemize}

\begin{tcolorbox}[breakable,colback=teal!4,colframe=teal!35!black,title={Before you move on}]
What does quamvīs mean? (\textquotedblleft{}However much; although (with the subjunctive)\textquotedblright{}.) Say it once more.
\end{tcolorbox}
